By [SGA 1, XIII 2.8], there is a specialization homomorphism (depending on
a path in $T$ from $\bar\eta$ to $o$, and hence defined up to conjugacy)
\begin{equation}\tag{6.1.3}
\pi_1^{\mathrm{mod}}(D_{\bar\eta}-D_{\bar\eta}\cap Z,\bar a)
\xrightarrow{\ \bar\eta\longrightarrow o\ }
\pi_1^{\mathrm{mod}}(D_o-D_o\cap Z,\bar a)
\end{equation}
such that the following assertions hold.

\subsubsection*{6.1.4.}
The local-monodromy homomorphisms (6.1.2.1)
\[
\prod_{\ell\ne\operatorname{char}(k)}\mathbb Z_\ell(1)
\longrightarrow
\pi_1^{\mathrm{mod}}(D_o-D_o\cap Z,\bar a)
\]
are obtained by composing the specialization arrow (6.1.3) with the local
monodromy homomorphisms for $D_{\bar\eta}$:
\begin{equation}\tag{6.1.4.1}
\prod_{\ell\ne\operatorname{char}(k)}\mathbb Z_\ell(1)
\longrightarrow
\pi_1^{\mathrm{mod}}(D_{\bar\eta}-D_{\bar\eta}\cap Z,\bar a)
\xrightarrow{\ \bar\eta\longrightarrow o\ }
\pi_1^{\mathrm{mod}}(D_o-D_o\cap Z,\bar a).
\end{equation}

\subsubsection*{6.1.5.}
The following diagram is commutative:
\begin{equation}\tag{6.1.5}
\begin{tikzcd}[column sep=huge,row sep=huge]
{} & \pi_1^{\mathrm{mod}}(\mathbb P^r-Z,\bar a) & {} \\
\pi_1^{\mathrm{mod}}(D_\eta-D_\eta\cap Z,\bar a)
  \arrow[ur,"i_{\eta*}"]
& {}
& \pi_1^{\mathrm{mod}}(D_o-D_o\cap Z,\bar a)
  \arrow[ul,"i_{o*}"'] \\
{} & \pi_1^{\mathrm{mod}}(D_{\bar\eta}-D_{\bar\eta}\cap Z,\bar a)
  \arrow[uu,"i_{\bar\eta*}"']
  \arrow[ul]
  \arrow[ur,"\bar\eta\longrightarrow o"'] & {} .
\end{tikzcd}
\end{equation}

Assume this for the moment, and first prove 6.1.1.  By the commutativity of
6.1.5, it is enough to prove that
\[
i_{\bar\eta*}:
\pi_1^{\mathrm{mod}}(D_{\bar\eta}-D_{\bar\eta}\cap Z,\bar a)
\longrightarrow
\pi_1^{\mathrm{mod}}(\mathbb P^r-Z,\bar a)
\]
is surjective.  For this, it is enough to prove that the analogous
homomorphism of the \underline{ordinary} fundamental groups
\begin{equation}\tag{6.1.6}
i_{\bar\eta*}:
\pi_1(D_{\bar\eta}-D_{\bar\eta}\cap Z,\bar a)
\longrightarrow
\pi_1(\mathbb P^r-Z,\bar a)
\end{equation}
is surjective.  In other words, for every étale covering
$E\to\mathbb P^r-Z$ with $E$ connected (and hence irreducible, since
$\mathbb P^r-Z$ is smooth and therefore unibranch [EGA IV, 18.10.3]), its
inverse image
\[
E_{\bar\eta}
=E\times_{\mathbb P^r-Z}(D_{\bar\eta}-D_{\bar\eta}\cap Z)
\]
over $D_{\bar\eta}-D_{\bar\eta}\cap Z$ must remain geometrically
irreducible.  It is enough to use the following variant of Bertini's theorem
([EGA V, citation left blank in the source]) for the composite
$E\to\mathbb P^r-Z\lhook\joinrel\longrightarrow\mathbb P^r$ and $s=1$.

\subsubsection*{6.1.6.1. Bertini's theorem.}
Let $k$ be a field, let $\mathbb P^r=\mathbb P^r_k$, and let
$f:X\to\mathbb P^r$ be a nonconstant morphism, where $X$ is a geometrically
irreducible algebraic scheme.  For every integer
$s<\dim\overline{f(X)}$, let $L_t$ denote the \underline{generic} linear
subspace of codimension $s$ in $\mathbb P^r$.  Then
$X_t=f^{-1}(L_t)$ is geometrically irreducible.
